\section{Implementation and On-Glasses Deployment}\label{sec:impl}

\subsection{Training}
\paragraph{Optimisation.}
\begin{itemize}
\item AdamW~\cite{loshchilov2019decoupled}, weight decay $10^{-4}$, batch size 32.
\item One-cycle schedule~\cite{smith2019super}: peak learning rate $3\times10^{-3}$ for the heads and $3\times10^{-4}$ for the backbone, 10\% warm-up, then cosine decay.
\item Up to 180 epochs, stopping early after 25 epochs without a lower validation loss.
\item Gradient clipping at 1.0, mixed precision, seed 42.
\end{itemize}
Training took 1\,h\,40\,min on one NVIDIA Tesla T4 (47\,s per epoch). Early stopping ended it after 125 epochs; the reported checkpoint is from epoch 100. Early stopping (patience 25 on the validation loss) ended the 180-epoch one-cycle schedule before it completed. Every run in this paper uses the same rule, so all comparisons are like for like; a completed schedule would, if anything, raise the absolute numbers.

\paragraph{Augmentation.} Besides the distractor pasting of Section~\ref{sec:ownership}:
\begin{itemize}
\item horizontal flip (no vertical flip, because a wearer's hand entering from the top would contradict the ownership prior);
\item shift, scale and rotation (${\pm}5\%$, ${\pm}10\%$, ${\pm}15^\circ$; probability 0.4);
\item perspective warp (0.2);
\item brightness and contrast (${\pm}20\%$; 0.5);
\item Gaussian blur (0.2);
\item mixup~\cite{zhang2018mixup} between images of the same scene, with $\lambda\sim\mathrm{Beta}(0.2,0.2)$ folded to $\geq0.5$ and the label of the dominant image.
\end{itemize}
All augmentation is on the fly; the dataset on disk is never modified. Appendix~\ref{app:aug} lists all parameters.

\paragraph{Export.} The model is exported to ONNX (opset 17, which bilinear grid sampling requires) and to FP16. The FP16 weights take 1.6\,MB.

\subsection{Deployment on Smart Glasses}\label{sec:device}
We deploy the 160\,px model on an engineering unit of a pair of smart glasses: a pre-production test device, not a retail product. Its hardware:
\begin{itemize}
\item a smart-glasses system-on-chip paired with a low-power microcontroller;
\item a single 12\,MP ultra-wide camera (105$^\circ$) in portrait orientation, with a dedicated low-power vision NPU;
\item Wi-Fi~6 and Bluetooth~5.4 to a companion phone;
\item a 160\,mAh battery, in a frame weighing under 45\,g.
\end{itemize}
Because it is an engineering device, its power management is not tuned for production, so absolute battery figures are pessimistic. We therefore compare power between operating modes on the same unit, rather than reporting battery life. We treat the unit as a measurement platform: its pre-production software stack is not public, so we report end-to-end device measurements rather than runtime-specific details.

\paragraph{Operating modes.} We compare two ways of running pointing localisation:
\begin{itemize}
\item \emph{On the glasses:} SpatialNet runs continuously on frames from the glasses' camera.
\item \emph{Offloading:} the glasses only capture frames and transmit them to the companion phone at 0.5 frames per second, the alternative in which localisation would run on the phone or in the cloud.
\end{itemize}

\paragraph{Measurement.}
\begin{itemize}
\item the inference latency per frame;
\item the peak additional RAM used during inference, measured over the idle application with the model already loaded (the 1.6\,MB of FP16 weights are stored separately);
\item the battery drain over one minute, in each operating mode, starting from the same charge level.
\end{itemize}

The camera delivers portrait frames, whereas most EgoPoint images are landscape. The model squashes both to a square, so the deployed input has a different aspect ratio from most training images. Real-frame behaviour is assessed in Section~\ref{sec:own}.
